% !TEX root = main.tex
\documentclass[12pt,a4paper]{ctexbook}
\usepackage{geometry}
\usepackage{tikz}
\usepackage{xcolor}
\usepackage{amsmath}
\usepackage{fancyhdr}
\usepackage{titlesec}
\usepackage{amsmath,amssymb,amsfonts}
\usepackage{geometry}
\usepackage{esint}
\usepackage{lastpage}
\usepackage{datetime2}
\usepackage{tcolorbox}
\usepackage{hyperref}
\usepackage{tikz-3dplot}
\usepackage{pgfplots}
\usetikzlibrary{calc}
\usetikzlibrary{matrix,3d}
\usetikzlibrary{tikzmark,backgrounds}
\usetikzlibrary{shapes.geometric}
\pgfdeclarelayer{background}
\pgfsetlayers{background,main}
\pgfplotsset{compat=1.13}
\DTMsetdatestyle{iso}
\setcounter{chapter}{-1}
\tcbuselibrary{theorems, skins, breakable}
\everymath{\displaystyle}
\geometry{margin=2.5cm}

% 定义统一的方框样式
\newtcbtheorem[number within=chapter]{definition}{定义}{%
    enhanced,
    colback=blue!5,      % 背景浅蓝色
    colframe=blue!75!black,
    coltitle=white,
    fonttitle=\bfseries,
    title={#2},          % 方框标题
    attach boxed title to top left={yshift=-2mm, xshift=5mm},
    boxed title style={sharp corners, rounded corners, colback=blue!75!black},
    separator sign none, % 不显示默认的冒号
    before skip=10pt,
    after skip=10pt,
}{def}

\newtcbtheorem[number within=chapter]{theorem}{定理}{%
    enhanced,
    colback=green!5,     % 背景浅绿色
    colframe=green!50!black,
    coltitle=white,
    fonttitle=\bfseries,
    title={#2},
    attach boxed title to top left={yshift=-2mm, xshift=5mm},
    boxed title style={sharp corners, rounded corners, colback=green!50!black},
    separator sign none,
    before skip=10pt,
    after skip=10pt,
    fontupper=\itshape
}{thm}

% 手动计数器，随章编号
\newcounter{myexample}[chapter]
\renewcommand{\themyexample}{\thechapter.\arabic{myexample}}

% 例题环境：椭圆标签在左，与文字同行
\newenvironment{myexample}
  {%
    \par\noindent                % 新段落，无缩进
    \refstepcounter{myexample}   % 计数器加1
    \tikz[baseline=(ell.base)]   % 对齐椭圆文字基线
      \node[ellipse, draw=black!40, fill=black!8,
            inner sep=1pt, font=\small,
            anchor=base] (ell) {例\themyexample};%
    \quad                        % 留一点间距
  }
  {%
    \par                         % 结束段落
  }


% 封面配色
\definecolor{titleblue}{HTML}{1B3A5C}
\definecolor{gold}{HTML}{C5A54E}

% 正文章节标题格式
\titleformat{\section}{\Large\bfseries\color{titleblue}\centering}{\thesection}{1em}{}
\titleformat{\subsection}{\large\bfseries\color{titleblue!80}}{\thesubsection}{1em}{}
